\section{Trasy i połączenia instalacyjne}
Każde połączenie ma dwa istniejące końce urządzenie:port, typ kabla, identyfikator trasy, pełną łamaną 3D w SITE, rzut, pion, zapas montażowy i długość zakupu. Dla odcinka o punktach $P_i=(x_i,y_i,z_i)$:
\[
L_r=\sum_i\sqrt{(x_{i+1}-x_i)^2+(y_{i+1}-y_i)^2},\quad
L_p=\sum_i|z_{i+1}-z_i|,
\]
\[
L_{\mathrm{inst}}=L_r+L_p+L_z,\qquad
L_{\mathrm{zak}}=\left\lceil L_{\mathrm{inst}}(1+f)\right\rceil.
\]
Dla kabla instalacyjnego $f=0{,}10$, a zapas końcowy wynosi zasadniczo 0,4--2 m, z większym zapasem na odcinku zewnętrznym; wartości są indywidualnie zapisane. Zaokrąglenie dotyczy każdego cięcia, więc suma zakupowa nie musi równać się zaokrąglonej długości całej grupy. Połączenia fabryczne oraz patchcordy ujęte w osprzęcie są oznaczone gwiazdką i nie są ponownie kupowane na metry.

\RouteTable

R1/R2 są trasami wspólnymi tylko mechanicznie: koryto ma rozdzielone funkcje IP, KD i SSWiN. SSP ma odrębne podpory i przejścia R3/R4; energia 230 V biegnie osobnym korytem. R5--R8 tworzą zewnętrzne drzewo kanalizacji teletechnicznej długości \DuctLength{} m, z zakupem \DuctPurchase{} m rury i trzema studzienkami poza pasem ruchu i obrysem otwarcia bramy. Rury nie są kablami; nie sumuje się ich do długości kabla. Współrzędne tras głównych określają korytarze prowadzenia, a indywidualne łamane kabli uwzględniają wszystkie odejścia i piony. To łamane kabli są podstawą długości i rysunku połączenia. JC03 przesunięto z $(65,11)$ na $(62,11)$ na istniejącym przebiegu R5; obrys 0,4 $\times$ 0,4 m kończy się na $x=62{,}2$, poza pasem od $x=63$. Nie zmienia to długości zasilania G01. Odgałęzienie C08 prowadzi teraz po R8 od $(63,11)$ do $(63,20)$; zaktualizowano przewód IP, PE i ich przedmiar.

Dane OSDP i zasilanie czytników stanowią osobne przewody i osobne połączenia. RS485 ma parę skręcaną 120 $\Omega$, zasilanie lokalne 2 $\times$ 0,75 mm$^2$, a długie zasilanie dwóch czytników G01 2 $\times$ 6 mm$^2$. Przy zacisku stosuje się krótkie przejście na przekrój dopuszczony przez producenta. Magistrala G01 prowadzi KD-05--czytnik IN--czytnik OUT, bez gwiazdy i długich odnóg; terminatory występują na obu końcach. Sygnał i wspólna masa nie są zastępowane ekranem.\cite{S-KD,S-READER}

Dla PIR przewód 6 $\times$ 0,5 mm$^2$ obejmuje oddzielną parę zasilania AUX i parę linii 2EOL; wolna para jest izolowana. Linie kontaktronów i przekaźników mają własne przewody sygnałowe. SSP wykorzystuje adresowalny przewód 1 $\times$ 2 $\times$ 1,0 mm, gdzie 1,0 mm oznacza średnicę żyły: przekrój wynosi 0,785 mm$^2$. Nie wolno w obliczeniu używać dla niego przekroju 1 mm$^2$.

Wybrano dedykowane zasilanie SSP kablem NHXH-J FE180/E90 z kompatybilnym systemem mocowań E90; pozostałe obwody 230 V mają odrębne YDY. Oznaczenie kabla nie dowodzi zachowania funkcji całego zespołu w pożarze. W rzeczywistej realizacji trzeba sprawdzić system podłoże--uchwyt--kabel i sposób uszczelnienia przejść. Nie deklaruje się przebadania ani certyfikacji instalacji.

\subsection*{Rejestr połączeń i długości}
Gwiazdka oznacza połączenie fabryczne albo długość już zawartą w pozycji patchcordów. Pełne współrzędne każdej łamanej podaje rejestr danych; tabela poniżej pozwala sprawdzić końce, podział długości i przedmiar.
\ConnectionTable
